% Kamp 1968, final substantive argument, printed pp. 112--117.

\subsection{Present progressive and the limits of ordinary past and future}
\label{sec:p09-kamp-present-progressive}

\paragraph{The formal question and its linguistic motivation.}
The thesis ends its substantive text by returning to the progressive
operation announced on pp.~18--19. Kamp considers two possible
formalizations and suggests that the use of progressive constructions in
natural language might reflect a conception of time as dense. That is
an interpretive suggestion, without an empirical demonstration in this
argument. The mathematical results below concern the explicitly defined
set operations and their expressibility; they do not establish a claim
about how speakers of any language represent time
\cite[pp.~18--19, 112--113]{Kamp1968}.

\begin{hbdefinition}{def:p09-kamp-progressive-language}{The present atom and the ordinary past/future language}
Fix a nonempty strict linear order \(\mathcal J=\langle T,<\rangle\).
For \(A\subseteq T\), put
\[
 \mathsf{PAST}(A)=\{x:\exists y<x\ (y\in A)\},\qquad
 \mathsf{FUTURE}(A)=\{x:\exists y>x\ (y\in A)\}.
\]
The source's \(\mathsf{TL}_0\) consists of finite formulas
generated from proposition letters by \(\neg,\land,P,F\),
where the unary symbols \(P,F\) denote these strict existential
operations. A one-input schema contains \(q_0\) and no other
proposition letter. The atom \(q_0\) denotes the input at the
current point, so the present is available throughout. Other finite
truth-functional connectives are Boolean abbreviations. There are
no named time points, additional predicate parameters, immediate-neighbor
operators, or isolation operators in this basis.

For \(a\leq b\) in \(T\), the symbols \((a,b)\) and
\([a,b]\) denote the open and closed intervals in the actual
order. They do not assert density or the existence of points between
the endpoints. For a formula \(\eta\) and input \(A\), write
\(\llbracket\eta\rrbracket_A\subseteq T\) for its truth set
\cite[pp.~18--19, 114]{Kamp1968}.
\end{hbdefinition}

\begin{hbdefinition}{def:p09-kamp-progressive-two-definitions}{Two distinct proposed progressive operations}
The initial proposal is
\[
\begin{aligned}
 \rho_0(t_0)={}&P_0(t_0)\land
   \exists t_1\bigl(t_1<t_0\land
        \forall t_2(t_1<t_2<t_0\to P_0(t_2))\bigr)\\
 &\land\exists t_3\bigl(t_0<t_3\land
        \forall t_4(t_0<t_4<t_3\to P_0(t_4))\bigr).
\end{aligned}
\]
Its operation is denoted \(\mathsf{PR}_0\). The revised
proposal additionally requires the two strict witnesses to satisfy
the input predicate:
\[
\begin{aligned}
 \rho(t_0)={}&P_0(t_0)\land
   \exists t_1\bigl(t_1<t_0\land P_0(t_1)\land
        \forall t_2(t_1<t_2<t_0\to P_0(t_2))\bigr)\\
 &\land\exists t_3\bigl(t_0<t_3\land P_0(t_3)\land
        \forall t_4(t_0<t_4<t_3\to P_0(t_4))\bigr).
\end{aligned}
\]
Write \(\mathsf{PR}\) for the operation defined by \(\rho\).
Both are equality-free first-order unary tenses with only \(t_0\)
free and only \(P_0\) in their unary signature. Their quantifiers
range over \(T\). In particular the revised operation is
\[
 \mathsf{PR}(A)=\{x\in T:\exists a<x<b\quad[a,b]\subseteq A\}.
\]
It is this revised operation, already specified in the initial
announcement on p.~19, which the final nonexpressibility result
addresses
\cite[pp.~19, 112--113]{Kamp1968}.
\end{hbdefinition}

\begin{hblemma}{thm:p09-kamp-progressive-definitions}{proved=annex:p9.kamp.progressive.definitions}{Dense, discrete and endpoint behavior of the two proposals}
On every strict linear order,
\(\mathsf{PR}(A)\subseteq\mathsf{PR}_0(A)\subseteq A\),
and both operations are false at a first or last time point.
On a dense order, \(\mathsf{PR}=\mathsf{PR}_0\), including
when the order has endpoints.

On a discrete order in the source's conditional-neighbor sense,
let \(T_*\) be the set of points having both an earlier and a
later point. Each \(x\in T_*\) has an immediate predecessor
\(x^-\) and successor \(x^+\), and
\[
 \mathsf{PR}_0(A)=A\cap T_*,\qquad
 \mathsf{PR}(A)=\{x\in T_*:x^-,x,x^+\in A\}.
\]
Thus the weak proposal equals the present operation on an endpoint-free
discrete order, whereas the revised one also tests the adjacent points.
On a discrete order with endpoints, the weak proposal has the expression
\[
 q_0\land P\neg(q_0\land\neg q_0)
       \land F\neg(q_0\land\neg q_0),
\]
which reduces to \(q_0\) when there are no endpoints.
For example, on \(\mathbb Z\) the singleton input \(\{0\}\)
has weak progressive value \(\{0\}\) and revised value
\(\varnothing\).

Using the source's strict, witness-first binary operators, on all
linear orders the revised operation has the expression
\[
                         q_0\land S(q_0,q_0)\land U(q_0,q_0).
\]
This expression belongs to \(\mathsf{TL}_1\), which has the
binary constructors \(S,U\); the nonexpressibility question below
concerns the smaller \(\mathsf{TL}_0\) basis
\cite[pp.~112--113]{Kamp1968}.
\end{hblemma}

\begin{hbdefinition}{def:p09-kamp-progressive-depths}{The two source directional depths}
Define even natural numbers \(d_-(\eta),d_+(\eta)\) by
\[
\begin{array}{c|cc}
 \eta&d_-(\eta)&d_+(\eta)\\ \hline
 q_0&0&0\\
 \neg\xi&d_-(\xi)&d_+(\xi)\\
 \xi\land\theta&\max(d_-(\xi),d_-(\theta))
                    &\max(d_+(\xi),d_+(\theta))\\
 P\xi&d_-(\xi)+2&d_+(\xi)\\
 F\xi&d_-(\xi)&d_+(\xi)+2
\end{array}
\]
The quantity \(d_-\) is the source's depth for its increasing
sequence argument, and \(d_+\) is its depth for the reversed
argument. In particular future adds no cost to \(d_-\), and
past adds no cost to \(d_+\). These are defined numerical
measures of finite syntax, not distances between time points
\cite[pp.~114--117]{Kamp1968}.
Unless a two-sided index set \(\mathbb Z\) is displayed,
all indices \(i,j,m,n\) below are nonnegative integers.
\end{hbdefinition}

\subsubsection{The source's sequence constructions with explicit local scope}

\begin{hblemma}{thm:p09-kamp-progressive-two-sided}{proved=annex:p9.kamp.progressive.two.sided}{The two-sided sequence construction}
Suppose a strictly increasing sequence \((a_n)_{n\in\mathbb Z}\)
of points of \(T\) is given. Put
\[
\begin{aligned}
 I&=\{x:\exists m,n\in\mathbb Z\ (a_m<x<a_n)\},\\
 A&=[a_{-2},a_2]\cup\{a_{2n}:n\in\mathbb Z\}.
\end{aligned}
\]
For every one-input \(\mathsf{TL}_0\)-formula \(\eta\),
\[
 \llbracket\eta\rrbracket_A\cap I
          \in\{\varnothing,I,A\cap I,I\setminus A\}.
\]
But \(\mathsf{PR}(A)=(a_{-2},a_2)\). It contains
\(a_0\) and excludes \(a_4\), although both belong to
\(A\cap I\). Thus \(\mathsf{PR}\) is not expressible
from the stated basis on any order with this sequence.
The sequence is not assumed cofinal in the whole order: the invariant
is on its convex hull \(I\), and temporal quantifiers still range
over all of \(T\)
\cite[pp.~113--114]{Kamp1968}.
\end{hblemma}

\begin{hblemma}{thm:p09-kamp-progressive-one-sided}{proved=annex:p9.kamp.progressive.one.sided}{Both one-sided sequence and depth constructions}
For a strictly increasing sequence \((a_n)_{n\geq0}\), put
\[
\begin{aligned}
 I_+&=\{x:a_0\leq x<a_n\text{ for some }n\geq0\},\\
 H_m^+&=\{x\in I_+:a_m<x\},\\
 A_i^+&=\{a_{2n}:n\geq0\}\cup[a_{i+2},a_{i+4}],
                   \qquad i\text{ even}.
\end{aligned}
\]
If \(i\geq d_-(\eta)\), put \(H=H_{d_-(\eta)}^+\).
Then
\[
 \llbracket\eta\rrbracket_{A_i^+}\cap H
             \in\{\varnothing,H,H\cap A_i^+,H\setminus A_i^+\}.
\]
The exact progressive image is
\(\mathsf{PR}(A_i^+)=(a_{i+2},a_{i+4})\).

For a strictly decreasing sequence \((b_n)_{n\geq0}\), instead put
\[
\begin{aligned}
 I_-&=\{x:b_n<x\leq b_0\text{ for some }n\geq0\},\\
 H_m^-&=\{x\in I_-:x<b_m\},\\
 A_i^-&=\{b_{2n}:n\geq0\}\cup[b_{i+4},b_{i+2}],
                   \qquad i\text{ even}.
\end{aligned}
\]
If \(i\geq d_+(\eta)\), put \(K=H_{d_+(\eta)}^-\).
Then
\[
 \llbracket\eta\rrbracket_{A_i^-}\cap K
             \in\{\varnothing,K,K\cap A_i^-,K\setminus A_i^-\},
\]
and
\(\mathsf{PR}(A_i^-)=(b_{i+4},b_{i+2})\).
Either sequence therefore supplies a counterinterpretation for any
proposed one-input definition. The inequalities on the even block
indices are part of the construction. A first or last point of the
whole order need not belong to either local tail
\cite[pp.~114--116]{Kamp1968}.
\end{hblemma}

\subsubsection{A finite version of the two-ended construction}

\begin{hbdefinition}{def:p09-kamp-progressive-window}{Separated marker chains and their filled intervals}
Choose a strictly increasing chain \((a_n)\) and a strictly decreasing chain
\((b_n)\), with \(a_n<b_m\) for every pair of available
indices. The indices may run over all nonnegative integers, as in
the source's two-chain construction, or over a finite segment
\(0,\ldots,N\). Choose even nonnegative \(i,j\); in the
finite case require \(\max(i,j)+6\leq N\). Put
\[
\begin{aligned}
 A_{ij}={}&\{a_{2n}:2n\text{ is available}\}
       \cup\{b_{2n}:2n\text{ is available}\}\\
       &\cup[a_{i+2},a_{i+4}]
        \cup[b_{j+4},b_{j+2}],\\
 C_{mn}={}&\{x\in T:a_m<x<b_n\}.
\end{aligned}
\]
The intervals include all their points in \(T\), not just the
selected markers. Both chains and their interval endpoints are
data for the interpretation; they are not constants available in
candidate formulas. The finite version is sufficient for the
unrestricted infinite-order conclusion.
\end{hbdefinition}

\begin{hblemma}{thm:p09-kamp-progressive-window}{proved=annex:p9.kamp.progressive.window}{The complete two-depth window invariant}
For the construction above, if
\(d_-(\eta)\leq i\) and \(d_+(\eta)\leq j\), put
\(C=C_{d_-(\eta),d_+(\eta)}\). Then
\[
 \llbracket\eta\rrbracket_{A_{ij}}\cap C
              \in\{\varnothing,C,C\cap A_{ij},C\setminus A_{ij}\}.
\]
This holds for both finite and infinite marker chains, with temporal
quantification over the whole order. Moreover,
\[
 \mathsf{PR}(A_{ij})
       =(a_{i+2},a_{i+4})\cup(b_{j+4},b_{j+2}).
\]
Both intervals lie inside the displayed central window. The points
\(a_{i+3},a_{i+6}\) both belong to the input and the window,
but only the first belongs to the progressive image
\cite[p.~117]{Kamp1968}.
\end{hblemma}

\begin{hbtheorem}{thm:p09-kamp-present-progressive}{proved=annex:p9.kamp.present.progressive}{Present progressive is not expressible by ordinary past and future on an infinite order}
For every infinite strict linear order \(\mathcal J=\langle T,<\rangle\),
there is no finite one-input \(\mathsf{TL}_0\)-schema \(\eta\)
such that
\[
                  \forall A\subseteq T\qquad
                      \llbracket\eta\rrbracket_A=\mathsf{PR}(A).
\]
The permitted primitive operations are exactly negation, conjunction,
strict existential \(\mathsf{PAST}\), and strict existential
\(\mathsf{FUTURE}\), with the present atom \(q_0\) available.
The order may have neither endpoint, just a first endpoint, just a
last endpoint, or both. No density, discreteness, completeness,
countable cofinality, or existence of separated infinite chains is
assumed in this conclusion.

For a candidate formula, it suffices to choose two separated finite
chains of length \(N+1\), where
\(N=\max(d_-(\eta),d_+(\eta))+6\), and use
\(A_{ij}\) with \(i=d_-(\eta)\), \(j=d_+(\eta)\).
Thus the general conclusion uses a construction from finitely many
selected points inside any infinite order
\cite[pp.~19, 112--117]{Kamp1968}.
\end{hbtheorem}

\paragraph{Relation to the complete-order expressive theorem.}
The revised progressive operation is first-order definable, and its
\(q_0\land S(q_0,q_0)\land U(q_0,q_0)\) representation works
on every linear order. The result above concerns the smaller ordinary
past/future language. It is consistent with, and illustrates the need
for a richer basis in, the earlier expressive-completeness result over
\(\mathsf{Com}\).

\begin{hbexample}{def:p09-kamp-progressive-boundaries}{Parity, endpoints and the bounded infinite order}
On \(\mathbb N\), let \(a_n=n\) and choose the even block
index \(i=0\). Then
\(A_0^+=\{0,2,4,6,\ldots\}\cup[2,4]\) has
\(\mathsf{PR}(A_0^+)=\{3\}\). Both three and six satisfy
the input, while only three satisfies the progressive operation.
If the block index were instead odd, the exact-image statement would
need additional qualifications: with \(i=1\), the same construction
contains every integer from two through six, and its progressive image
is \(\{3,4,5\}\).

At the first point of \(\mathbb N\), both progressive proposals
are false even when the input contains that point. On the bounded
infinite order \(\omega+1\), a decreasing sequence indexed by all
nonnegative integers does not exist, although both a first and a last
point exist. The finite-window construction still applies, by choosing
the required markers among finitely many natural numbers below the
last point. The annex verifies these boundary examples.
\end{hbexample}
